\begin{example}
    Sea $k\in \mathbb{Z}$ tal que $k\geq1$, sea $a \neq 0$ y sea $\tau < 0$. 
    Entonces, la sucesión $(an^k\log^\tau n)_{n\geq2}$ es equidistribuida módulo $1$.
\end{example}

\begin{resolve}
    Sea $x_n = an^k\log^\tau n$ para cada $n \geq 2$. Definimos $g(x) = a(x+1)^k\log^\tau (x+1)$
    en~$[1,\infty)$, entonces $g(n) = x_{n+1}$ para cada $n \geq 1$.

    La derivada $k$-ésima de $g(x)$ es
    \[
        g^{(k)}(x) = a \log^{\tau-k}(x+1)\left(P(\log(x+1))\right)
    \]
    donde $P$ es un polinomio de grado $k$ con coeficiente principal $k!$.

    Debido a que $\tau < 0$, entonces 
    \[
        \lim_{x \to \infty} g^{(k)}(x) = 0
        \qquad \text{y} \qquad
        \lim_{x \to \infty} x \left|g^{(k)}(x)\right| = \infty
    \]

    Para comprobar la monotonía de $g^{(k)}$, calculamos su derivada,
    \[
        g^{(k+1)}(x) = a \,\frac{\log^{\tau-k-1}(x+1)}{(x+1)}\left(Q(\log(x+1))\right)
    \]
    donde $Q$ es un polinomio de grado $k+1$ con coeficiente principal $\tau k!$.

    Dado que $\frac{\log^{\tau-k-1}(x+1)}{(x+1)} > 0$ para cada $x \geq 1$ y como $\tau k! < 0$,
    \[
        \lim_{x \to \infty} Q(x) = -\infty
    \] 
    entonces existe $x_0 \geq 1$ tal que $Q(\log(x+1)) < 0$ para cada $x \geq x_0$. 
    En consecuencia, $g^{(k+1)}(x)$ es de signo constante para cada $x \geq x_0$, y 
    por tanto, $g^{(k)}(x)$ es monótona para cada $x \geq x_0$.

    Por el teorema \ref{teo:fejer_cap4}, $(g(n))$ es equidistribuida módulo $1$ y 
    por lo tanto, $(x_n)$ también lo es.
\end{resolve}